\scanpage{63}
\mspage{31}
\uline{Conjecture B'} Soit $K$ un corps de type fini sur $\mathbb{Q}$ (ou sur un corps cyclotomique, cela suffira peut-être), $\bar K$ une clôture algébrique de $K$, $\mathcal U$ une courbe de Poincaré sur $\bar K$, de complétée $\widehat{\mathcal U}=\mathcal X$. Considérons une action de $\Gamma=\Gamma_{\bar K/K}$ sur $\mathcal U$, compatible avec son action sur $\bar K$, d'où une action de $\Gamma$ sur $\mathcal X$. Ceci posé :

il existe un point fixe et un seul de $\Gamma^\circ$ agissant sur $\mathcal X$

($\Gamma^\circ$, noyau du caractère cyclotomique $\Gamma\xrightarrow{\chi}\widehat{\mathbb Z}^{\!*}$).

La différence avec la conjecture B, que celle-ci renforce, c'est qu'ici on ne suppose \uline{pas} que l'action de $\Gamma$ sur $\mathcal U$ normalise l'action d'un groupe profini $\pi$, agissant librement sur $\mathcal U$ de façon \unsure{que|q.} $\mathcal U/\pi$ soit une courbe alg. anabélienne sur $\bar K$.
\margin{62}
